\subsection{Where the Democratic Dividend Actually Comes From}
\label{sec:12.1.1}
AI that flags disinformation, exposes corruption, protects an election's integrity, or helps a citizen engage with government: these are real applications, and none of them is what makes AI protect democratic values rather than undermine them.

Freedom House's 2023 assessment of internet freedom found at least twenty-two countries where law already requires or incentivizes a platform to deploy machine learning against content the state disfavors, almost never described in those terms: the operative language is national security or a threat of disinformation \autocite{funk2023freedom}. India's IT Rules, amended in February 2026 to require platforms to label AI-generated content and act on certain takedown notices within three hours, were defended as protection against synthetic disinformation and fraud — the application the list above would have praised on its own terms. A lawyer who has litigated India's free-speech cases documents that enforcement concentrates on parody and criticism of the prime minister rather than on the harms the rules name \autocite{gupta2026three}.

The application did not fail at what it was built to do. It worked as built, inside a system missing the independent court or press whose absence is what made the result a liability, and whether those exist is a fact about a country's politics settled before any AI system is built.
